\chapter*{Giới thiệu về báo cáo}
\addcontentsline{toc}{chapter}{Giới thiệu về báo cáo}

Môn học \textbf{Tính toán song song} cung cấp các kiến thức cơ bản và nâng cao về lập trình đa xử lý trên hệ thống máy tính hiện đại. Trong báo cáo này, nhóm thực hiện lập trình và thử nghiệm các bài toán số học cơ bản bằng công cụ \textbf{OpenMP} (Open Multi-Processing) - một API hỗ trợ lập trình song song bộ nhớ chia sẻ phổ biến và mạnh mẽ trên C/C++ và Fortran.

Nội dung báo cáo tập trung vào hai bài tập thực hành chính:
\begin{enumerate}
    \item \textbf{Bài tập 1: Tính tích vô hướng của hai vector}: Thực hiện song song hóa vòng lặp tính toán tích vô hướng trên vector, thử nghiệm hiệu năng và tính đúng đắn với các kích thước dữ liệu khác nhau.
    \item \textbf{Bài tập 2: Giải hệ phương trình tuyến tính}: Sử dụng hai phương pháp phổ biến là phương pháp khử Gauss và phương pháp lặp Jacobi. Báo cáo này sẽ trình bày chi tiết về phần thuật toán và lập trình song song của phương pháp khử Gauss. Phần phương pháp Jacobi sẽ do thành viên khác trong nhóm phụ trách hoàn thiện.
\end{enumerate}

Báo cáo bao gồm phân tích lý thuyết thuật toán, giải thích chi tiết từng đoạn mã nguồn C++ có tích hợp OpenMP và đưa ra các kết quả thử nghiệm thực tế từ chương trình.

